\uuid{vF8T}
\titre{Changements de variables}
\auteur{hanani ; mzari}
\organisation{exo7}
\licence{CC BY-NC-SA 4.0}
\datecreate{2020-01-20}
\niveau{L1}
\module{Analyse}
\chapitre{Calcul d'intégrales}
\sousChapitre{Autre}
\difficulte{3}
\theme{Calculs d'intégrales, changements de variables}
\exo7qcmid{lille-2.integrales.017}

\contenu{
\begin{question}
\qclassification{127.99}
\qtags{motcle=changements de variables}
On pose $\displaystyle I=\int _{0}^{1/\sqrt{2}}\frac{x\ \mathrm{d}x}{\sqrt{1-x^2}}$. Parmi les affirmations suivantes, cocher celles qui sont vraies :
\begin{answers}
\bad{Le changement de variable $t=1-x^2$ donne $\displaystyle I=-\int _0^{1/\sqrt{2}}\frac{\mathrm{d}t}{2\sqrt{t}}$.}
\good{Le changement de variable $t=1-x^2$ donne $\displaystyle I=\int _{1/2}^1\frac{\mathrm{d}t}{2\sqrt{t}}$.}
\bad{Le changement de variable $t=x^2$ donne $\displaystyle I=\int _0^{1/\sqrt{2}}\frac{\mathrm{d}t}{2\sqrt{1-t}}$.}
\good{Le changement de variable $x=\sin t$ donne $\displaystyle I=\int _{0}^{\pi/4}\sin t\, \mathrm{d}t$.}
\end{answers}
\vskip3mm
\begin{explanations}
Avec $t=1-x^2$, on a : $t(0)=1$, $t(1/\sqrt{2})=1/2$, $\mathrm{d}t=-2x\mathrm{d}x$ et $\displaystyle I=\int _{1/2}^1\frac{\mathrm{d}t}{2\sqrt{t}}$.
\vskip0mm
Avec $t=x^2$, on a : $t(0)=0$, $t(1/\sqrt{2})=1/2$, $\mathrm{d}t=2x\mathrm{d}x$ et $\displaystyle I=\int _0^{1/2}\frac{\mathrm{d}t}{2\sqrt{1-t}}$.
\vskip0mm
Avec $x=\sin t$, on a : $t(0)=0$, $t(1/\sqrt{2})=\pi/4$, $\sqrt{1-x^2}=\cos t$, $\mathrm{d}x=\cos t\, \mathrm{d}t$ et
$$\displaystyle I=\int _0^{\pi/4}\sin t\mathrm{d}t.$$
\end{explanations}
\end{question}
}
